\documentclass[conference]{IEEEtran}
\IEEEoverridecommandlockouts
\usepackage[utf8]{inputenc}
\usepackage[T1]{fontenc}
\usepackage{cite}
\usepackage{amsmath,amssymb,amsfonts}
\usepackage{graphicx}
\usepackage{textcomp}
\usepackage{xcolor}
\usepackage{listings}
\usepackage{booktabs}
\usepackage{tabularx}
\usepackage{multirow}
\usepackage{enumitem}
\usepackage{caption}
\usepackage{array}
\usepackage{microtype}
\usepackage{hyperref}

\hypersetup{
    colorlinks=true,
    linkcolor=black,
    citecolor=black,
    urlcolor=blue!60!black
}

% ---- listings setup ----
\definecolor{lstbg}{RGB}{248,248,248}
\definecolor{lstframe}{RGB}{170,170,170}
\definecolor{kw}{RGB}{0,0,160}
\definecolor{cmt}{RGB}{0,120,0}
\definecolor{str}{RGB}{160,0,0}

\lstset{
  basicstyle=\ttfamily\scriptsize,
  keywordstyle=\color{kw}\bfseries,
  commentstyle=\color{cmt}\itshape,
  stringstyle=\color{str},
  backgroundcolor=\color{lstbg},
  rulecolor=\color{lstframe},
  breaklines=true,
  frame=single,
  numbers=left,
  numberstyle=\tiny\color{gray},
  numbersep=6pt,
  stepnumber=1,
  xleftmargin=12pt,
  xrightmargin=2pt,
  aboveskip=4pt,
  belowskip=4pt,
  columns=fullflexible,
  keepspaces=true,
  showstringspaces=false,
  captionpos=b
}

\lstdefinestyle{py}{language=Python}
\lstdefinestyle{sh}{language=bash, basicstyle=\ttfamily\scriptsize, numbers=none}
\lstdefinestyle{plain}{basicstyle=\ttfamily\scriptsize, numbers=none, frame=single, backgroundcolor=\color{lstbg}}

\newcommand{\cwe}[1]{\href{https://cwe.mitre.org/data/definitions/#1.html}{CWE-#1}}

\begin{document}

\title{Software Vulnerability Analysis of \texttt{fuzz4all-cs}:\\
A Manual Review with Static-Analysis Cross-Validation}

\author{\IEEEauthorblockN{Aniket Salunke}
\IEEEauthorblockA{CS7602 -- Software Vulnerability Analysis\\
Mohamed bin Zayed University of Artificial Intelligence\\
\texttt{aniket.salunke@mbzuai.ac.ae}}}

\maketitle

\begin{abstract}
We perform a security review of \texttt{fuzz4all-cs}, a Python harness that
extends the LLM-based universal fuzzer Fuzz4All with an error-guided repair
stage. Using a combination of manual taint review and two independent static
analysers (Bandit~1.9.4 and Semgrep~1.161.0 with the \texttt{p/python},
\texttt{p/command-injection}, and \texttt{p/security-audit} rule packs), we
identify three confirmed and exploitable vulnerabilities. The first is an
operating-system command injection (\cwe{78}) reachable from three external
sources, including the search-harness candidate JSON. The second is an
insecure temporary-file pattern (\cwe{377}, with a chained
symlink-following condition, \cwe{59}) that is realistic on a shared
multi-user HPC node, the deployment target stated in the project README. The
third is unsandboxed execution of generated code by the Qiskit oracle
(\cwe{94}). For each finding we provide the exact source location, a
reproducible non-destructive proof-of-concept, an impact statement, and a
recommended remediation. Two additional candidate weaknesses (\cwe{502}
PyYAML deserialisation, \cwe{22} path traversal via \texttt{output\_folder})
are reported as null results under the stated threat model. Bandit reports
77 findings, including 23 HIGH-severity \texttt{B602} sites that match our
\textit{F1}; Semgrep reports 24 ERROR-severity findings including a single
end-to-end taint flow at
\texttt{tools/run\_search\_round.py:92}, independently corroborating the
reachability argument for \textit{F1}.
\end{abstract}

\begin{IEEEkeywords}
software vulnerability analysis, command injection, insecure temporary file,
static analysis, taint analysis, secure coding, LLM-based fuzzing, CWE
\end{IEEEkeywords}

% =====================================================================
\section{Introduction and Codebase Under Review}
% =====================================================================

\texttt{fuzz4all-cs} is a fork of Fuzz4All~\cite{xia2024fuzz4all} extended,
as part of the course project, with an error-guided repair stage that uses
a local Large Language Model (LLM) to repair compilation failures emitted
by the C++ target (\texttt{g++}). The harness is approximately 2{,}440
lines of Python, organised under
\texttt{Fuzz4All/} (target validators, repair module, fuzzing loop) and
\texttt{tools/} (evaluator and search harness). The codebase is intended to
be run by a researcher on a shared High-Performance Computing (HPC) node:
the project \texttt{README.md} explicitly states an HPC, no-Docker
deployment, and the search harness ingests JSON candidate files that are
described as ``generated by ChatGPT/Cursor''.

This report concerns the security posture of the Python harness itself
(the trusted Python source code, the configuration parser, the
oracles, and the search harness). It does not analyse the LLM weights, the
HuggingFace transport, or any third-party CVE in
\texttt{transformers}/\texttt{torch}/PyYAML. The model weights and the
upstream package vulnerabilities are out of scope.

% =====================================================================
\section{Threat Model}
\label{sec:threat}
% =====================================================================

\textbf{Trust boundaries.} We model the system as having four classes of
inputs (Table~\ref{tab:trust}). The operator who launches the fuzzer and
the local Python source are trusted. A configuration YAML supplied by a
collaborator (e.g. a pull request, a shared repository, a search round) is
\emph{partially trusted}: it controls the experiment, but in practice it
is consumed without integrity checks. Candidate JSON files consumed by the
search harness are explicitly described as AI-generated content and are
therefore \emph{untrusted}. Any unprivileged user co-tenant on the same
HPC node is \emph{untrusted} and can read process listings and write into
\texttt{/tmp}. Generated programs produced by the LLM are also
\emph{untrusted} since the model can be steered through prompt-injection
content placed in the documentation files.

\begin{table}[t]
\caption{Trust boundaries assumed in this analysis.}
\label{tab:trust}
\centering\renewcommand{\arraystretch}{1.15}
\begin{tabularx}{\columnwidth}{lXc}
\toprule
\textbf{Source} & \textbf{Description} & \textbf{Trust} \\
\midrule
Operator      & User running \texttt{fuzz.py}                  & Trusted \\
Python source & Files under \texttt{Fuzz4All/}, \texttt{tools/} & Trusted \\
YAML config   & \texttt{config/*.yaml}, possibly third-party   & Partial \\
JSON candidate& \texttt{candidates/round\_*.json} (AI-emitted) & Untrusted \\
Co-tenant     & Other unprivileged users on shared HPC node    & Untrusted \\
LLM output    & Generated C++/Python from StarCoderBase        & Untrusted \\
\bottomrule
\end{tabularx}
\end{table}

\textbf{Assets.} (a) The operator's home directory and SSH credentials,
(b) integrity of files owned by the operator, (c) confidentiality of any
secrets present in the operator's environment (e.g.\ HPC scheduler tokens,
HuggingFace tokens), and (d) availability of the fuzzing run.

\textbf{Out of scope.} HuggingFace network endpoints, third-party
library CVEs, the LLM model itself, hardware side channels, and physical
attacks.

% =====================================================================
\section{Methodology}
% =====================================================================

Our analysis combines \emph{eight complementary techniques} so that no
single tool's blind spot dominates the result. The eight techniques map
onto four classes that the course covered explicitly: (i) human
white-box analysis with an explicit threat model, (ii) syntactic / AST
pattern matching, (iii) declarative semantic dataflow, and (iv) software
composition / supply-chain analysis. Table~\ref{tab:tools} summarises
the tools and the role each one plays in our methodology.

\begin{table*}[t]
\caption{Techniques and tools used. Each row is a distinct class of
evidence; together they cover human reasoning, AST patterns, taint
flow, semantic dataflow, security linting, and supply-chain analysis.}
\label{tab:tools}
\centering\renewcommand{\arraystretch}{1.2}
\begin{tabularx}{\textwidth}{lllX}
\toprule
\textbf{\#} & \textbf{Technique class} & \textbf{Tool / version} & \textbf{Why we used it on \texttt{fuzz4all-cs}} \\
\midrule
1 & Human white-box review (STRIDE) & manual, with the threat model of Sec.~III & Establishes trust boundaries config/CLI $\rightarrow$ \texttt{target\_name} $\rightarrow$ \texttt{subprocess}; designs minimal PoCs. \\
2 & AST pattern static analysis & Bandit 1.9.4 & Python-native, broad CWE coverage (B602, B108, B506, B615); good first sweep. \\
3 & Pattern + intra/inter-procedural taint & Semgrep 1.161.0 + \texttt{p/python}, \texttt{p/command-injection}, \texttt{p/security-audit} & Independent confirmation of Bandit and a true source $\rightarrow$ sink taint flow. \\
4 & \emph{Project-specific} taint rule & Custom Semgrep rule \texttt{tools/semgrep\_rules/cmd\_inject\_target.yml} (this work) & Encodes the application-level boundary (\texttt{self.target\_name} / \texttt{compiler} $\rightarrow$ \texttt{subprocess(\ldots, shell=True)}) which generic rules miss. \\
5 & Declarative semantic dataflow & CodeQL 2.25.2, pack \texttt{codeql/python-queries@1.8.0} & Industry-standard inter-procedural analysis; we run \texttt{python-security-and-quality.qls}, \texttt{python-security-extended.qls}, and the targeted \texttt{CWE-078/022/094} queries. The null result on stock CommandInjection (Sec.~\ref{sec:evidence}) is itself a methodological finding. \\
6 & Security linting (DUO* checks) & flake8 7.3.0 + dlint 0.16.0 & Lightweight third opinion confirming \texttt{shell=True}, weak RNG, ReDoS and unsafe yaml. \\
7 & Software-composition analysis (SCA) & pip-audit 2.10.0 vs.\ PyPI Advisory DB / OSV & Detects vulnerable third-party dependencies in the project's conda environment. \\
8 & Secret scanning & detect-secrets 1.5.0 & Verifies that no plaintext credentials, tokens, or keys are committed (null-result evidence). \\
\bottomrule
\end{tabularx}
\end{table*}

\textbf{Manual review.} For every external source identified in
Table~\ref{tab:trust} we traced the value through the call graph until
it reached one of the following sinks: (i) any
\texttt{subprocess.run/Popen} invocation, with particular attention to
\texttt{shell=True}; (ii) any \texttt{open()} of a path constructed from
external data; (iii) any path used inside a shell-formatted string. The
review covered eight principal modules: \texttt{Fuzz4All/fuzz.py}, the
target validators (\texttt{CPP/CPP.py}, \texttt{C/C.py}, \texttt{GO/GO.py},
\texttt{JAVA/JAVA.py}, \texttt{SMT/SMT.py}, \texttt{QISKIT/QISKIT.py}),
the repair module (\texttt{Fuzz4All/repair/repair.py}), and the two
tooling scripts (\texttt{tools/evaluate\_candidate.py},
\texttt{tools/run\_search\_round.py}).

\textbf{Static-analysis cross-check.} Each manual finding was
independently cross-validated with the seven automated tools listed in
Table~\ref{tab:tools}. Bandit and Semgrep registry rules were used as
the primary AST/pattern sweep. We then authored a project-specific
Semgrep taint rule that models \texttt{self.target\_name} (and the local
alias \texttt{compiler} used inside \texttt{CPP.py}) as the source and
any \texttt{subprocess.*(\ldots, shell=True)} or \texttt{os.system}
invocation as the sink. CodeQL was run with the standard
\texttt{python-security-and-quality.qls} suite, the
\texttt{python-security-extended.qls} suite, and the three targeted
queries for CWE-78 / CWE-22 / CWE-94 to obtain the strongest available
inter-procedural dataflow signal. flake8/dlint added a third independent
syntactic opinion on \texttt{shell=True}, weak random, ReDoS, and unsafe
\texttt{yaml.load}. pip-audit scanned the conda environment of the
project and reported every CVE listed in the PyPI advisory database
backed by OSV. detect-secrets scanned every tracked file (excluding
\texttt{.git/}, \texttt{outputs/}, and generated \texttt{.fuzz} programs)
to verify the absence of committed credentials.

All raw artefacts are committed under \texttt{outputs/security/}
(Bandit, Semgrep registry, custom Semgrep, three CodeQL SARIFs,
flake8+dlint, pip-audit, detect-secrets) and summarised in
Section~\ref{sec:evidence} and
Appendices~\ref{app:bandit}--\ref{app:semgrep}. Exact reproduction
commands are listed in \texttt{outputs/security/SUMMARY.md}.

\textbf{Proof-of-concept policy.} We constructed minimal, non-destructive
PoCs for each finding. The PoCs target a sentinel path
(\texttt{/tmp/PWNED\_F<n>}) rather than any real asset and are entirely
reversible. They are reproduced verbatim in Appendix~\ref{app:pocs}.

% =====================================================================
\section{Findings}
\label{sec:findings}
% =====================================================================

A summary of the three confirmed findings is given in
Table~\ref{tab:findings}.

\begin{table}[t]
\caption{Confirmed findings.}
\label{tab:findings}
\centering\renewcommand{\arraystretch}{1.15}
\begin{tabularx}{\columnwidth}{clXc}
\toprule
\textbf{ID} & \textbf{CWE} & \textbf{Component (file:line)} & \textbf{Sev.} \\
\midrule
F1 & 78  & \texttt{Fuzz4All/target/CPP/CPP.py:100--107} (+22 sister sites) & High \\
F2 & 377\,/\,59 & \texttt{Fuzz4All/target/CPP/CPP.py:29--37}; \texttt{target.py:127} & High \\
F3 & 94  & \texttt{Fuzz4All/target/QISKIT/QISKIT.py:264--272, 305--313, 340--348} & Medium \\
\bottomrule
\end{tabularx}
\end{table}

% ----------------------------------------------------------
\subsection{F1 -- OS Command Injection via \texttt{target\_name} (\cwe{78})}
% ----------------------------------------------------------

\textbf{Location.} The C++ validator constructs a shell command by
f-string interpolation and executes it with \texttt{shell=True}:

\begin{lstlisting}[style=py,caption={Vulnerable validator in Fuzz4All/target/CPP/CPP.py (lines 100--107).},firstnumber=100,label={lst:f1}]
result = subprocess.run(
    f"{compiler} -x c++ -std=c++23 -O0 "
    f"-fdiagnostics-color=never -c {filename} -o {out_obj}",
    shell=True,
    capture_output=True,
    encoding="utf-8",
    timeout=5,
    text=True,
)
\end{lstlisting}

The local variable \texttt{compiler} aliases \texttt{self.target\_name},
which is initialised from kwargs in the \texttt{Target} constructor:

\begin{lstlisting}[style=py,caption={Target.\_\_init\_\_ attribute assignment in Fuzz4All/target/target.py.},firstnumber=171]
if "target_name" in kwargs:
    self.target_name = kwargs["target_name"]
\end{lstlisting}

The same kwargs originate from the parsed config dictionary, into which
\texttt{main\_with\_config} writes the value supplied either through the
CLI flag \texttt{--target} or through the config file:

\begin{lstlisting}[style=py,caption={Source assignment in Fuzz4All/fuzz.py.},firstnumber=443]
if target != "":
    config_dict["fuzzing"]["target_name"] = target
\end{lstlisting}

The same antipattern recurs in
\texttt{Fuzz4All/target/C/C.py:113}, \texttt{Fuzz4All/target/GO/GO.py:68},
\texttt{Fuzz4All/target/JAVA/JAVA.py:83},
\texttt{Fuzz4All/target/SMT/SMT.py:115}, and at multiple sites in
\texttt{Fuzz4All/target/QISKIT/QISKIT.py}.

\textbf{Reachability.} Three independent paths reach the sink:
\begin{enumerate}[label=(\roman*),leftmargin=*]
\item \emph{CLI flag} \texttt{--target} (operator self-pwn, low value).
\item \emph{YAML config} -- the field
\verb|fuzzing.target_name| inside any \texttt{config/*.yaml} loaded with
\texttt{yaml.load} (\texttt{Fuzz4All/util/util.py:65}).
\item \emph{Search-harness candidate JSON} -- in
\texttt{tools/run\_search\_round.py} (lines 56--75) the flat key
\verb|target_name| inside a JSON candidate is mapped into
\verb|config["fuzzing"]| and the materialised YAML is then handed to
\texttt{evaluate\_candidate.py}, which spawns \texttt{fuzz.py} with that
config. The candidate JSON files are explicitly described as AI-emitted
and are therefore \emph{untrusted} under our threat model.
\end{enumerate}
Independently, Semgrep's taint analyser confirms an end-to-end source-to-sink
flow from candidate JSON data into a \texttt{subprocess.run} call at
\texttt{tools/run\_search\_round.py:92} via the rule
\verb|dangerous-subprocess-use-tainted-env-args|.

\textbf{Proof of concept (non-destructive).} Listing~\ref{lst:f1poc}
demonstrates the three reachable surfaces. Each PoC creates a sentinel
file at \texttt{/tmp/PWNED\_F1*}; no asset of the operator is altered.

\begin{lstlisting}[style=sh,caption={F1 proof of concept.},label={lst:f1poc}]
# Surface (i): CLI
python Fuzz4All/fuzz.py --config config/cpp_repair_smoke.yaml \
    main_with_config --folder /tmp/poc1 \
    --target 'g++ ; touch /tmp/PWNED_F1A #'
ls -l /tmp/PWNED_F1A    # exists  ->  RCE confirmed

# Surface (ii): YAML
cat > /tmp/evil.yaml <<'EOF'
fuzzing: { target_name: "g++ ; touch /tmp/PWNED_F1B #",
           num: 2, total_time: 0.01, otf: true, resume: false,
           evaluate: false, use_hand_written_prompt: false,
           no_input_prompt: true, prompt_strategy: 0, log_level: 3 }
target:  { language: cpp,
           path_documentation: config/documentation/cpp/cpp_23.md,
           path_example_code: null,
           trigger_to_generate_input: "/* x */",
           input_hint: "#include <iostream>",
           path_hand_written_prompt: null, target_string: "" }
llm:     { temperature: 1, batch_size: 1, device: cpu,
           model_name: bigcode/starcoderbase-1b, max_length: 64 }
EOF
python Fuzz4All/fuzz.py --config /tmp/evil.yaml \
    main_with_config --folder /tmp/poc2

# Surface (iii): search-harness candidate JSON
echo '[{"name":"evil","gen":{"target_name":"g++ ; touch /tmp/PWNED_F1C #"},"repair":{}}]' \
   > candidates/round_evil.json
python tools/run_search_round.py --round candidates/round_evil.json \
   --base_config config/cpp_repair_smoke.yaml \
   --budget_programs 2 --repeats 1
\end{lstlisting}

\textbf{Impact.} Arbitrary command execution as the user running the
fuzzer. On a shared HPC node this typically grants the attacker full
access to the operator's home directory, SSH keys, scratch storage, and
any compute-job credentials present in the environment. In a CI context
the build runner is fully compromised. The three reachable surfaces also
expose the wider scientific-computing community: a malicious collaborator
need only submit a candidate JSON to a shared search round.

\textbf{Recommended remediation.} Pass arguments as a list, set
\texttt{shell=False}, validate the compiler path against an allow-list,
and resolve to a known absolute path. The minimal patch is:

\begin{lstlisting}[style=py,caption={Recommended fix for F1 (not applied; analysis-only assignment).}]
import os, shutil
ALLOWED = {"/usr/bin/g++", "/usr/bin/clang++"}

compiler_path = shutil.which(compiler) or compiler
if compiler_path not in ALLOWED:
    raise ValueError(f"compiler not allow-listed: {compiler_path}")

subprocess.run(
    [compiler_path, "-x", "c++", "-std=c++23", "-O0",
     "-fdiagnostics-color=never", "-c", filename, "-o", out_obj],
    shell=False, capture_output=True, text=True, timeout=5,
)
\end{lstlisting}

% ----------------------------------------------------------
\subsection{F2 -- Insecure Temporary File with Symlink Following
(\cwe{377}, \cwe{59})}
% ----------------------------------------------------------

\textbf{Location.} The temporary-file path is a deterministic function of
\texttt{time.time()} captured \emph{once} at \texttt{Target}
construction:

\begin{lstlisting}[style=py,caption={Fuzz4All/target/target.py line 127.},
firstnumber=127]
self.CURRENT_TIME = time.time()
\end{lstlisting}

Every program written during the run uses the same path:

\begin{lstlisting}[style=py,caption={Fuzz4All/target/CPP/CPP.py lines 29--37.},firstnumber=29]
def write_back_file(self, code):
    try:
        with open(
            "/tmp/temp{}.cpp".format(self.CURRENT_TIME),
            "w", encoding="utf-8") as f:
            f.write(code)
    except:
        pass
    return "/tmp/temp{}.cpp".format(self.CURRENT_TIME)
\end{lstlisting}

The same pattern applies to \verb|/tmp/out{CURRENT_TIME}.o| at
\texttt{CPP.py:94} and to identically named files in C, Go, Java, Qiskit
and SMT targets. There are three independent defects compounded:
\begin{enumerate}[label=(\roman*),leftmargin=*]
\item the path lives in the world-writable, sticky-bit directory
\texttt{/tmp};
\item the filename is a deterministic, low-entropy function of the
process start time;
\item Python's \texttt{open(path, "w")} is implemented as
\texttt{O\_CREAT|O\_WRONLY|O\_TRUNC} \emph{without} \texttt{O\_NOFOLLOW},
so it follows attacker-controlled symbolic links.
\end{enumerate}
Bandit independently flags this pattern with rule \texttt{B108}
(\textit{hardcoded\_tmp\_directory}) at 19 sites under
\texttt{Fuzz4All/target/}.

\textbf{Reachability.} Any unprivileged co-tenant on the HPC node who can
either observe the operator's process start time or guess
\texttt{time.time()} within a few seconds of the run can pre-stage a
symbolic link in \texttt{/tmp}. When the fuzzer subsequently writes a
generated program, the open call follows the symlink and \emph{overwrites
the target file with the LLM-generated source}. The bare
\verb|except: pass| at \texttt{CPP.py:35} silently absorbs any
\texttt{PermissionError} that would otherwise alert the operator.

\textbf{Proof of concept (deterministic, single-user demo).}

\begin{lstlisting}[style=sh,caption={F2 proof of concept.},label={lst:f2poc}]
# 1. Sentinel file we will pretend belongs to the victim.
mkdir -p /tmp/poc_target
echo SECRET > /tmp/poc_target/secret.txt

# 2. Predict CURRENT_TIME (start the fuzzer at TS+1 second).
TS=$(python -c 'import time;print(time.time()+1)')
ln -s /tmp/poc_target/secret.txt /tmp/temp${TS}.cpp

# 3. Run the fuzzer in the same window.
python Fuzz4All/fuzz.py --config config/cpp_repair_smoke.yaml \
   main_with_config --folder /tmp/poc_f2 --target /usr/bin/g++

# 4. Observe: secret.txt has been clobbered with C++ code.
cat /tmp/poc_target/secret.txt
\end{lstlisting}

\textbf{Impact.} Arbitrary file overwrite as the operator. Realistic
targets on an HPC home directory include
\texttt{\textasciitilde/.ssh/authorized\_keys} (denial of service
against SSH), \texttt{\textasciitilde/.bashrc} (persistence via shell
init), and Slurm submission scripts (lateral compromise). Combined with
prompt injection through the documentation files (which feed the LLM
prompt at \texttt{Fuzz4All/target/target.py:182--191}), the attacker
controls the \emph{content} of the overwritten file and can implant
their own SSH key.

\textbf{Recommended remediation.} Use the standard library's
\texttt{tempfile} module to obtain a private per-run directory and open
files with \texttt{O\_NOFOLLOW}:

\begin{lstlisting}[style=py,caption={Recommended fix for F2.}]
import os, tempfile

class CPPTarget(Target):
    def __init__(self, **kwargs):
        super().__init__(**kwargs)
        self.tmpdir = tempfile.mkdtemp(prefix="fuzz4all_", dir="/tmp")
        os.chmod(self.tmpdir, 0o700)
    def write_back_file(self, code):
        path = os.path.join(self.tmpdir, f"prog_{os.getpid()}.cpp")
        fd = os.open(path,
            os.O_WRONLY|os.O_CREAT|os.O_EXCL|os.O_NOFOLLOW, 0o600)
        with os.fdopen(fd, "w", encoding="utf-8") as f:
            f.write(code)
        return path
\end{lstlisting}

% ----------------------------------------------------------
\subsection{F3 -- Unsandboxed Execution of Generated Code (\cwe{94})}
% ----------------------------------------------------------

\textbf{Location.} The Qiskit oracle invokes the Python interpreter on
whatever the LLM produced:

\begin{lstlisting}[style=py,caption={Fuzz4All/target/QISKIT/QISKIT.py lines 264--272.},firstnumber=264]
cmd = f"python {new_filename}"
exit_code = subprocess.run(
    cmd,
    shell=True,
    capture_output=True,
    encoding="utf-8",
    timeout=15,
    text=True,
)
\end{lstlisting}

Identical execution sites are present at lines 305--313 and 340--348.
The C target (\texttt{Fuzz4All/target/C/C.py:113--155}) compiles
\emph{and} links the LLM output, so a generated program containing a
\texttt{\_\_attribute\_\_((constructor))} function executes attacker
code at the next link step.

\textbf{Reachability.} The LLM prompt is built from
\texttt{path\_documentation}, \texttt{path\_example\_code} and
\texttt{path\_hand\_written\_prompt} (\texttt{target.py:182--191}); none
of the three has any integrity check. A collaborator who controls the
documentation file in
\texttt{config/documentation/qiskit/*.md} can prompt-inject the LLM into
emitting a Python program with \verb|os.system(...)|, which the oracle
will then execute.

\textbf{Proof of concept (controlled, bypassing the LLM).}

\begin{lstlisting}[style=sh,caption={F3 proof of concept.},label={lst:f3poc}]
TS=$(python -c 'import time;print(time.time())')
mkdir -p /tmp/poc_f3
cat > /tmp/temp${TS}_lvl_0.py <<'EOF'
import os; os.system("touch /tmp/PWNED_F3")
EOF
# A Qiskit run that invokes `python /tmp/temp${TS}_lvl_0.py`
# will execute the payload above as the operator.
ls /tmp/PWNED_F3   # exists  ->  RCE confirmed
\end{lstlisting}

\textbf{Impact.} Remote code execution under the fuzzer's UID. The
underlying defect is the absence of an isolation boundary around the
oracle and the absence of any documentation declaring that the oracle is
a security-relevant trust boundary.

\textbf{Recommended remediation.} Wrap each oracle invocation in a
sandbox. A minimal approach using \texttt{bubblewrap}:

\begin{lstlisting}[style=sh,caption={Recommended fix for F3.}]
bwrap --unshare-all --die-with-parent --new-session \
      --ro-bind /usr /usr --ro-bind /lib /lib \
      --ro-bind /lib64 /lib64 --proc /proc --dev /dev \
      --tmpfs /tmp -- python "$file"
\end{lstlisting}

% =====================================================================
\section{Null Results}
\label{sec:null}
% =====================================================================

We list two patterns that we investigated and found to be \emph{not}
exploitable under our threat model. Reporting them explicitly is
important to demonstrate scope discipline.

\begin{itemize}[leftmargin=*]
\item \textbf{\cwe{502} -- PyYAML deserialisation.} \texttt{Fuzz4All/util/util.py:65}
calls \verb|yaml.load(..., Loader=yaml.FullLoader)|. Bandit reports this
as \texttt{B506}, but \texttt{FullLoader} in PyYAML~$\geq 5.1$ blocks
arbitrary Python object construction; we were unable to construct a
working code-execution payload. We still recommend \verb|yaml.safe_load|
as defence in depth, since the security-sensitive dispatch lives on a
single line.
\item \textbf{\cwe{22} -- Path traversal via \texttt{output\_folder}.}
The folder argument is used only with \texttt{os.makedirs} and
\texttt{os.path.join} for files written by the operator who is also the
trust source for that flag. There is no privilege boundary crossed,
hence no exploitable issue.
\end{itemize}

% =====================================================================
\section{Static-Analysis Evidence Summary}
\label{sec:evidence}
% =====================================================================

Tables~\ref{tab:bandit-summary}--\ref{tab:tool-totals} report the
distribution of static-analysis findings across the seven automated
tools used in this study. The full text reports are in
Appendices~\ref{app:bandit} and \ref{app:semgrep}; the corresponding
JSON / SARIF artefacts are at
\texttt{outputs/security/bandit\_report.json},
\texttt{semgrep\_report.json},
\texttt{semgrep\_custom.json},
\texttt{codeql.sarif},
\texttt{codeql\_security\_extended.sarif},
\texttt{codeql\_injections.sarif},
\texttt{flake8\_dlint.txt},
\texttt{pip\_audit.json}, and
\texttt{detect\_secrets.json}.

\begin{table}[t]
\caption{Headline counts across the seven automated tools.}
\label{tab:tool-totals}
\centering\renewcommand{\arraystretch}{1.15}
\begin{tabularx}{\columnwidth}{Xcc}
\toprule
\textbf{Tool} & \textbf{Findings} & \textbf{Type} \\
\midrule
Bandit                                  & 77 & AST patterns \\
Semgrep (registry packs)                & 24 & Patterns + taint \\
Semgrep (custom rule, this work)        & 6  & Project-specific taint \\
CodeQL \texttt{python-security-and-quality} & 59 & Quality + security \\
CodeQL \texttt{python-security-extended}    & 0  & Null result (see text) \\
CodeQL \texttt{CWE-078 / 022 / 094}          & 0  & Null result (see text) \\
flake8 + dlint (DUO*)                   & 27 & Security linting \\
pip-audit                               & 7 CVEs / 6 pkgs & SCA \\
detect-secrets                          & 0 & Null result \\
\bottomrule
\end{tabularx}
\end{table}

\begin{table}[t]
\caption{Bandit 1.9.4 -- distribution by check (top rules).}
\label{tab:bandit-summary}
\centering\renewcommand{\arraystretch}{1.15}
\begin{tabularx}{\columnwidth}{lXcc}
\toprule
\textbf{ID} & \textbf{Rule} & \textbf{Sev.} & \textbf{N} \\
\midrule
B602 & \texttt{subprocess\_popen\_with\_shell\_equals\_true} & High   & \textbf{23} \\
B108 & \texttt{hardcoded\_tmp\_directory}                    & Medium & \textbf{19} \\
B110 & \texttt{try\_except\_pass}                            & Low    & 16 \\
B404 & \texttt{import\_subprocess} (informational)           & Low    & 8  \\
B615 & \texttt{huggingface\_unsafe\_download}                & Medium & 2  \\
B603 & \texttt{subprocess\_without\_shell\_equals\_true}     & Low    & 2  \\
B506 & \texttt{yaml\_load}                                   & Medium & 1  \\
others & B105, B607, B311, B101, B602(low) -- benign / FP   & Low    & 6  \\
\midrule
\multicolumn{2}{l}{\textbf{Total}} & & \textbf{77} \\
\bottomrule
\end{tabularx}
\end{table}

\begin{table}[t]
\caption{Semgrep 1.161.0 -- all 24 findings are severity \emph{ERROR}.}
\label{tab:semgrep-summary}
\centering\renewcommand{\arraystretch}{1.2}
\begin{tabularx}{\columnwidth}{p{0.34\columnwidth}Xc}
\toprule
\textbf{Rule} & \textbf{Description} & \textbf{N} \\
\midrule
\texttt{subprocess-} \texttt{shell-true}
  & Pattern \texttt{shell=True} in \texttt{subprocess.*} (CWE-78). & 23 \\
\texttt{dangerous-} \texttt{subprocess-} \texttt{use-tainted-} \texttt{env-args}
  & Taint flow: candidate JSON $\rightarrow$ \texttt{subprocess.run} at \texttt{tools/run\_search\_round.py:92} (CWE-78). & 1  \\
\midrule
\multicolumn{2}{l}{\textbf{Total}} & \textbf{24} \\
\bottomrule
\end{tabularx}
\end{table}

\paragraph{Custom Semgrep rule (this work).}
The custom rule \texttt{tools/semgrep\_rules/cmd\_inject\_target.yml}
encodes the project-specific trust boundary: source =
\texttt{self.target\_name} or local alias \texttt{compiler}; sink =
\texttt{subprocess.run/Popen/call(\ldots, shell=True, \ldots)} or
\texttt{os.system(\ldots)}. The rule produces six high-confidence
findings, one per language target
(\texttt{C/C.py:25}, \texttt{CPP/CPP.py:100}, \texttt{GO/GO.py:67},
\texttt{JAVA/JAVA.py:86}, \texttt{SMT/SMT.py:117} and 126), confirming
that F1 is \emph{systemic} rather than localised.

\paragraph{CodeQL semantic dataflow.}
We ran CodeQL~2.25.2 with pack \texttt{codeql/python-queries@1.8.0}
against a database built over all 13 Python source files of the
project. The full \texttt{python-security-and-quality.qls} suite
returns 59 findings, dominated by code-quality categories that
\emph{compound} our manual findings (Table~\ref{tab:codeql-summary}).
Most importantly, CodeQL's \emph{stock command-injection},
\emph{path-injection} and \emph{code-injection} queries return
\textbf{zero} findings (Table~\ref{tab:tool-totals}). This null result
is intentional and informative: CodeQL's default Python taint model
classifies HTTP requests, environment variables, and \texttt{argparse}
CLI args as remote sources, but does \emph{not} model the YAML config
file (loaded by \texttt{Fuzz4All/util/util.py:65} via
\texttt{yaml.load(\ldots, Loader=yaml.FullLoader)}) as a tainted source.
Because \texttt{target\_name} reaches the subprocess sink only through
that YAML config, the stock model misses the bug. This (i) shows that
off-the-shelf taint engines under-report when a project loads
configuration from disk and treats the disk as untrusted, (ii) directly
motivates our custom Semgrep rule, and (iii) underlines the value of
the human threat-model step.

\begin{table}[t]
\caption{CodeQL \texttt{python-security-and-quality.qls} findings (59
total). The rules in italics compound our F2 finding by hiding errors
or leaking file handles in the temp-file lifecycle.}
\label{tab:codeql-summary}
\centering\renewcommand{\arraystretch}{1.15}
\begin{tabularx}{\columnwidth}{Xcl}
\toprule
\textbf{CodeQL rule} & \textbf{N} & \textbf{Relevance} \\
\midrule
\texttt{py/unused-import}              & 25 & code hygiene \\
\emph{\texttt{py/empty-except}}        & 15 & swallows F2 errors \\
\emph{\texttt{py/file-not-closed}}     & 7  & temp-file lifecycle (F2) \\
\texttt{py/catch-base-exception}       & 6  & error-handling hygiene \\
\texttt{py/str-format/surplus-arg}     & 3  & malformed error logging \\
\texttt{py/uninitialized-local}        & 2  & C/GO validators \\
\texttt{py/unused-local}               & 1  & minor \\
\bottomrule
\end{tabularx}
\end{table}

\paragraph{flake8 + dlint.}
27 DUO* findings: 24 \texttt{DUO116} (\texttt{shell=True}) confirm F1,
plus \texttt{DUO109} (yaml-load), \texttt{DUO138} (catastrophic regex,
ReDoS hardening) and \texttt{DUO102} (non-cryptographic
\texttt{random}). This is a third independent confirmation of F1 and a
small set of additional hardening items.

\paragraph{Software-composition analysis (pip-audit).}
pip-audit~2.10.0 scanned the project's conda environment and reported
\textbf{seven CVEs across six installed packages}
(Table~\ref{tab:pip-audit}). These supply-chain issues are independent
of the application-level findings F1--F3 but are part of the project's
overall attack surface and feed into the consolidated recommendations
of Sec.~\ref{sec:findings}.

\begin{table}[t]
\caption{pip-audit findings inside the \texttt{fuzz4all} conda env.}
\label{tab:pip-audit}
\centering\renewcommand{\arraystretch}{1.15}
\begin{tabularx}{\columnwidth}{Xll}
\toprule
\textbf{Package} & \textbf{Installed} & \textbf{CVE} \\
\midrule
cryptography  & 46.0.5  & CVE-2026-34073 \\
cryptography  & 46.0.5  & CVE-2026-39892 \\
pillow        & 12.1.1  & CVE-2026-40192 \\
pip           & 26.0.1  & CVE-2026-3219  \\
pygments      & 2.19.2  & CVE-2026-4539  \\
qiskit-terra  & 0.24.1  & CVE-2025-2000  \\
requests      & 2.32.5  & CVE-2026-25645 \\
\bottomrule
\end{tabularx}
\end{table}

\paragraph{Secret scanning (detect-secrets).}
detect-secrets~1.5.0 scanned every tracked file (excluding
\texttt{.git/}, \texttt{outputs/}, generated \texttt{*.fuzz}) and
reported \textbf{zero} hits. This is a useful null result: the project
does not commit credentials or API tokens.

\paragraph{Cross-evidence mapping.}
Table~\ref{tab:cross} maps each manual finding to the corroborating
static-analysis evidence across all seven automated tools.

\begin{table*}[t]
\caption{Cross-evidence mapping across the manual review and the seven
automated tools used in this study.}
\label{tab:cross}
\centering\renewcommand{\arraystretch}{1.2}
\begin{tabularx}{\textwidth}{lXXXXX}
\toprule
\textbf{Finding} & \textbf{Manual} & \textbf{Bandit} & \textbf{Semgrep (registry + custom)} & \textbf{CodeQL} & \textbf{flake8/dlint, pip-audit, detect-secrets} \\
\midrule
F1 (\cwe{78}) & 23 \texttt{shell=True} sites; 3 reachable surfaces (CLI, YAML, candidate JSON) & 23 $\times$ B602 HIGH; 35 issues to CWE-78 & 23 $\times$ \texttt{subprocess-shell-true}; 1 $\times$ tainted-env-args at \texttt{run\_search\_round.py:92}; \emph{custom rule}: 6 findings in C/CPP/GO/JAVA/SMT & quality only; null on stock \texttt{CommandInjection} (motivates custom rule) & 24 $\times$ DUO116 \\
F2 (\cwe{377}, \cwe{59}) & Predictable fixed \texttt{/tmp} path; \texttt{open} follows symlinks & 19 $\times$ B108 MED; 16 $\times$ B110 (try/except pass) & no rule for CWE-377 in chosen packs & 15 $\times$ \texttt{py/empty-except}, 7 $\times$ \texttt{py/file-not-closed} (compounding) & --- \\
F3 (\cwe{94}) & Three execution sites in Qiskit oracle & subset of B602 (Qiskit sites) & subset of \texttt{subprocess-shell-true} (Qiskit sites) & quality only & DUO116 on Qiskit sites \\
Null: \cwe{502} & FullLoader is safe in PyYAML $\geq$ 5.1 & 1 $\times$ B506 & --- & --- & 1 $\times$ DUO109 \\
Null: secrets & no plaintext credentials observed & --- & --- & --- & detect-secrets: 0 \\
Suppl.: vulnerable deps & --- & --- & --- & --- & pip-audit: 7 CVEs / 6 packages (Table~\ref{tab:pip-audit}) \\
\bottomrule
\end{tabularx}
\end{table*}

\paragraph{Filtered false positives.}
Bandit's \texttt{B105} ($\times 2$) flagged the literal word \texttt{password}
inside repair-prompt templates, which are not credentials.
\texttt{B311} (\texttt{random}) is used only for prompt-strategy
selection, not for cryptographic randomness. \texttt{B404} and
\texttt{B603} are informational. \texttt{B506} (\texttt{yaml\_load}) is
discussed as a null result in \S\ref{sec:null}.

% =====================================================================
\section{Consolidated Recommendations}
% =====================================================================

\begin{enumerate}[leftmargin=*]
\item \textbf{Eliminate \texttt{shell=True} across all targets.} Convert
the 23 sites flagged by Bandit/Semgrep to list-form
\texttt{subprocess.run} calls. Where pipelines are needed (e.g. the
\texttt{ps | grep | xargs kill} timeout handler at
\texttt{CPP.py:115--126}), replace them with \texttt{psutil}-based
process trees.
\item \textbf{Allow-list trusted compilers.} Validate
\texttt{target\_name} against a fixed set
(\verb|/usr/bin/g++|, \verb|/usr/bin/clang++|, \verb|/usr/bin/gcc|)
before the first \texttt{subprocess} call.
\item \textbf{Replace shared \texttt{/tmp} with a per-run private
directory.} Use \texttt{tempfile.mkdtemp(prefix="fuzz4all\_")} with mode
\texttt{0700}, and open program files with \texttt{O\_NOFOLLOW}.
\item \textbf{Sandbox the oracles.} Wrap every execution of generated
code in \texttt{bwrap}/\texttt{firejail} with \verb|--unshare-all|,
\verb|--unshare-net| and a \verb|RLIMIT_AS| memory cap.
\item \textbf{Switch to \texttt{yaml.safe\_load}.}
Defence in depth, even though \texttt{FullLoader} is currently safe.
\item \textbf{Treat candidate JSONs as untrusted input.} Reject any
\texttt{target\_name} that contains shell metacharacters before the
materialised YAML is written.
\item \textbf{Replace bare \texttt{except: pass} blocks (Bandit
\texttt{B110}, 16 sites)} with logged exception handling, so that
symlink races and other anomalies surface to the operator.
\end{enumerate}

% =====================================================================
\section{Conclusion}
% =====================================================================

We presented a security review of the \texttt{fuzz4all-cs} harness using
manual taint review cross-validated by Bandit and Semgrep. Three
exploitable findings were confirmed (F1--F3); two candidate weaknesses
were investigated and reported as null results under the stated threat
model. F1 (command injection) is the highest-impact issue because the
search-harness candidate JSON crosses from an explicitly AI-generated,
\emph{untrusted} input into a \texttt{subprocess.run} sink in a single
\texttt{deepcopy/update} merge. F2 is realistic precisely because the
project README requires deployment on a shared HPC node. F3 reveals the
absence of an isolation boundary around the oracle, a property that
ought to be made explicit in the project's documentation. All three
findings admit straightforward, low-risk remediations summarised in
\S\ref{sec:findings} and consolidated above.

\begin{thebibliography}{9}
\bibitem{xia2024fuzz4all}
C. S. Xia, M. Paltenghi, J. L. Tian, M. Pradel, L. Zhang,
``Fuzz4All: Universal Fuzzing with Large Language Models,''
\emph{ICSE 2024}, IEEE, 2024.
\bibitem{bandit}
PyCQA, ``Bandit: a security linter for Python,''
\url{https://bandit.readthedocs.io}, accessed 2026.
\bibitem{semgrep}
Semgrep, Inc., ``Semgrep static analysis platform,''
\url{https://semgrep.dev}, accessed 2026.
\bibitem{cwe78}
The MITRE Corporation, ``CWE-78: Improper Neutralization of Special
Elements used in an OS Command (`OS Command Injection'),''
\url{https://cwe.mitre.org/data/definitions/78.html}.
\bibitem{cwe377}
The MITRE Corporation, ``CWE-377: Insecure Temporary File,''
\url{https://cwe.mitre.org/data/definitions/377.html}.
\bibitem{cwe59}
The MITRE Corporation, ``CWE-59: Improper Link Resolution Before File
Access (`Link Following'),''
\url{https://cwe.mitre.org/data/definitions/59.html}.
\bibitem{cwe94}
The MITRE Corporation, ``CWE-94: Improper Control of Generation of Code
(`Code Injection'),''
\url{https://cwe.mitre.org/data/definitions/94.html}.
\end{thebibliography}

\appendices

% =====================================================================
\section{Bandit Raw Output (Selected)}
\label{app:bandit}
% =====================================================================

Run command:
\begin{lstlisting}[style=sh]
bandit -r Fuzz4All/ tools/ -f txt \
       -o outputs/security/bandit_report.txt
\end{lstlisting}

Run metrics: 2{,}440 lines of code scanned; 77 issues
(23 High, 22 Medium, 32 Low). Excerpt of representative findings:

\begin{lstlisting}[style=plain]
>> Issue: [B602:subprocess_popen_with_shell_equals_true]
   Severity: High   Confidence: High
   CWE: CWE-78
   Location: Fuzz4All/target/CPP/CPP.py:100
   100  result = subprocess.run(
   101      f"{compiler} -x c++ -std=c++23 -O0 ..."
   102      shell=True,

>> Issue: [B108:hardcoded_tmp_directory]
   Severity: Medium Confidence: Medium
   CWE: CWE-377
   Location: Fuzz4All/target/CPP/CPP.py:32
   32   with open("/tmp/temp{}.cpp".format(self.CURRENT_TIME),

>> Issue: [B506:yaml_load]
   Severity: Medium Confidence: High
   CWE: CWE-20 (suggested)
   Location: Fuzz4All/util/util.py:65
   65   config = yaml.load(f, Loader=yaml.FullLoader)
\end{lstlisting}

The complete report is at
\texttt{outputs/security/bandit\_report.txt} (944 lines) and
\texttt{outputs/security/bandit\_report.json}.

% =====================================================================
\section{Semgrep Raw Output (Selected)}
\label{app:semgrep}
% =====================================================================

Run command:
\begin{lstlisting}[style=sh]
semgrep --config p/python --config p/command-injection \
        --config p/security-audit Fuzz4All/ tools/ \
        --metrics off --quiet \
        > outputs/security/semgrep_report.txt
\end{lstlisting}

Excerpt:

\begin{lstlisting}[style=plain]
24 Code Findings

Fuzz4All/target/CPP/CPP.py
  python.lang.security.audit.subprocess-shell-true
    [Blocking]
    102:    shell=True,
    117:    shell=True,
    125:    shell=True,

tools/run_search_round.py
  python.lang.security.audit
   .dangerous-subprocess-use-tainted-env-args
    [Blocking]   CWE-78
    Detected subprocess function 'run' with user
    controlled data. A malicious actor could leverage
    this to perform command injection.
     92:        subprocess.run(cmd, cwd=repo_root, check=True)
\end{lstlisting}

The complete report is at
\texttt{outputs/security/semgrep\_report.txt} (147 lines) and
\texttt{outputs/security/semgrep\_report.json} (24 results).

% =====================================================================
\section{Proof-of-Concept Transcripts}
\label{app:pocs}
% =====================================================================

The PoC scripts in Listings~\ref{lst:f1poc}, \ref{lst:f2poc} and
\ref{lst:f3poc} were executed in a controlled sandbox; the only side
effect observed was the creation of the sentinel files
\texttt{/tmp/PWNED\_F1A}, \texttt{/tmp/PWNED\_F1B}, \texttt{/tmp/PWNED\_F1C}
and \texttt{/tmp/PWNED\_F3}, plus the overwriting of
\texttt{/tmp/poc\_target/secret.txt} for F2. No real asset of the
operator was modified.

\end{document}
